\chapter{逆関数}
    \section{連続関数の逆関数は連続}
        \begin{proof}
            \quad\par
            まず元の関数の定義域が開区間の場合について示す。
            $f$を区間$I = (a_1,a_2)$上の連続関数とし、$J \coloneq f(I) =: (b_1,b_2)$とする。
            $f$には逆関数が存在するものとする(必然的に$f$は狭義単調増加/減少なものに限られる)。
            $f$が狭義単調増加な場合について示すが、狭義単調減少の場合も同様にして示せる。
            \par
            $a\in I$を任意に選ぶ。
            $f(a)$に於ける$f^{-1}$の右連続性を否定して矛盾が生じることを示す。
            右連続性を否定すると、ある$\alpha>0$が存在して、任意に小さい如何なる$\gamma>0$に対してもある適当な$y\in[f(a),f(a)+\gamma]$が存在して、$f^{-1}(y) \geq a+\alpha$となる。
            \par
            $a+\alpha$が$f$の定義域に含まれることを確認しておく。
            $f^{-1}(y)$なる値が存在していることと逆関数の定義から$f^{-1}(y)$は$f$の定義域に含まれていなくてはならない。
            よって$I \ni a < a+\alpha\leq f^{-1}(y) \in I$であり、$f$の定義域$I$は区間であると冒頭で仮定しているから$a+\alpha\in I$でなくてはならない。
            \par
            先述の通り$f^{-1}(y) \geq a+\alpha$であって、$f$は単調増加だから$y \geq f(a+\alpha)$である。
            これと$y \leq f(a)+\gamma$より$f(a+\alpha) \leq f(a)+\gamma$であり、$\gamma$は任意に小さくできるから$\gamma\to 0$とすると$f(a+\alpha) \leq f(a)$となり、$f$が狭義単調増加であるという前提に矛盾する。
            \par
            同様に$f(a)$に於ける$f^{-1}$の左連続性を否定しても矛盾が生じることが示される。
            よって$f^{-1}$は$f(a)$で連続であり、$a$は$I$内の任意の点だから$f(a)$は$J$内の任意の点となり得るため$f^{-1}$は$J$上連続である。
            \par
            次に$f$の定義域が閉区間$I=[a_1,a_2]$と表される場合について示す。
            半開区間の場合の証明は開区間の場合と閉区間の場合の証明を組み合わせて直ちに得られる
            (半開区間を隣接する閉区間と開区間に分けて考える。境界となる点での連続性は開区間の場合で証明済みである)。
            $J \coloneq f(I) =: [b_1,b_2]$とする。
            開区間での右連続性の証明と同様にして$b_1$での右連続性が示される。
            また開区間での左連続性の証明と同様にして$b_2$での左連続性が示される。
            以上より$f^{-1}$は定義域の端点を含めて連続である。
            \begin{figure}[H]
                \centering
                \includegraphics[keepaspectratio, scale=0.7]
                {\currfiledir/figs/continuity_of_inverse_function/graph.pdf}
                \caption{$a,\alpha,\gamma$の様子}
            \end{figure}
        \end{proof}
    \section{可微分な関数\texorpdfstring{$f$}{f}の逆関数は\texorpdfstring{$f'\left(f^{-1}(x)\right)\neq 0$}{f'(f\string^(-1)(x)) ≠ 0}なる点\texorpdfstring{$x$}{x}で微分可能であり、\texorpdfstring{$\deriv{f^{-1}}{x}{}(x) = 1/f'\left(f^{-1}(x)\right)$}{f\string^(-1)'(x) = 1/f'(f\string^(-1)(x))}。}
        \begin{proof}
            \quad\par
            $f$を開区間$I$上の可微分な関数とし、$f$には逆関数が存在するとする(必然的に$f$は狭義単調増加/減少なものに限られる)。
            まず$I$が開区間である場合について示す。
            $f$が狭義単調増加な場合について示すが、狭義単調減少の場合も同様にして示せる。
            $b \in \dom{f^{-1}}$を任意にとる。
            $k>0$を$b+1 \in \dom{f^{-1}}$となるように十分小さくとる。
            $f$の微分可能性より$k = f'\left(f^{-1}(b)\right)\left(f^{-1}(b+k) - f^{-1}(b)\right) + o\left(f^{-1}(b+k) - f^{-1}(b)\right)$であるから
            \begin{align*}
                \frac{f^{-1}(b+k) - f^{-1}(b)}{k} &= \frac{f^{-1}(b+k) - f^{-1}(b)}{f'\left(f^{-1}(b)\right)\left(f^{-1}(b+k) - f^{-1}(b)\right) + o\left(f^{-1}(b+k) - f^{-1}(b)\right)} \\
                &= \frac{1}{f'\left(f^{-1}(b)\right) + \frac{o\left(f^{-1}(b+k) - f^{-1}(b)\right)}{f^{-1}(b+k) - f^{-1}(b)}}
            \end{align*}
            $k\to 0$とすると$f^{-1}$の連続性から$f^{-1}(b+k) - f^{-1}(b) \to 0$となるので
            \[ \lim_{k\to 0} \frac{f^{-1}(b+k) - f^{-1}(b)}{k} = \frac{1}{f'\left(f^{-1}(b)\right)} \]
            \begin{figure}[H]
                \centering
                \includegraphics[keepaspectratio, scale=0.6]
                {\currfiledir/figs/differentiability_of_inverse_function/graph.pdf}
                \caption{様子}
            \end{figure}
        \end{proof}